\documentclass[10pt,letterpaper]{article}
\usepackage{cornell-notes}

\title{Chapter 8: Multiple Processor Systems}
\author{}
\date{}

\setDocTitle{Chapter 8: Multiple Processor Systems}
\setDocSubtitle{Operating Systems}
\setDocVersion{v1.0}
\setDocStatus{Study Companion}
\setDocOwner{Operating Systems Cornell Notes}
\setDocAuthor{}
\setDocDate{\today}
\setDocSummary{Cornell-format study and review notes for Chapter 8: Multiple Processor Systems.}

\setCornellCollection{Operating Systems Cornell Notes}
\setCornellUnitType{Chapter}
\setCornellUnitNumber{8}
\setCornellUnitTitle{Multiple Processor Systems}
\setCornellCompanionLabel{Study and review companion}

\hypersetup{pdftitle={Chapter 8: Multiple Processor Systems},pdfsubject={Cornell notes},pdfauthor={}}

\begin{document}
\maketitle
\section*{Learning Objectives}
\begin{studybox}{By the end of this chapter, you should be able to}
\begin{enumerate}
  \item Compare shared-memory multiprocessors, multicomputers, virtual machines, and distributed systems.
  \item Explain UMA, NUMA, cache coherence, and common interconnection structures.
  \item Compare separate, master-slave, and symmetric multiprocessor operating-system organizations.
  \item Analyze synchronization and scheduling choices for multiprocessors.
  \item Trace low-level and user-level message transmission in a multicomputer.
  \item Explain RPC, distributed shared memory, multicomputer scheduling, and load-balancing policy.
  \item Compare document-, file-, object-, and coordination-based middleware.
\end{enumerate}
\end{studybox}

\begin{studybox}{Section Map}
\hyperlink{sec-8-1}{8.1} \quad \hyperlink{sec-8-2}{8.2} \quad \hyperlink{sec-8-3}{8.3} \quad \hyperlink{sec-8-4}{8.4} \quad \hyperlink{sec-8-5}{8.5}
\end{studybox}

\section*{Cornell Cue-and-Notes Area}
\small
\begin{longtable}{C N}
\rowcolor{navy}
\color{white}\textbf{Cue / Recall Prompt} &
\color{white}\textbf{Notes}\\
\endfirsthead
\rowcolor{navy}
\color{white}\textbf{Cue / Recall Prompt} &
\color{white}\textbf{Notes (continued)}\\
\endhead
\rowcolor{ice}\multicolumn{2}{l}{\hypertarget{sec-8-1}{}\textbf{Section 8.1: Multiprocessors}}\\*\hline
\textbf{Shared-memory model} & Two or more processors access a common memory abstraction. Hardware interconnection, cache coherence, and kernel synchronization determine whether performance scales.\\\hline
\textbf{8.1.1 Hardware} & Bus-based UMA systems offer uniform access but limited bandwidth; crossbars and multistage networks scale connections; NUMA systems distribute memory so placement affects latency. Private caches require a coherence protocol.\\\hline
\textbf{8.1.2 OS organizations} & Each CPU can run a separate system, one master can handle kernel work for slaves, or a symmetric multiprocessor can let every CPU enter one shared kernel. SMP balances flexibility against synchronization complexity.\\\hline
\textbf{8.1.3 Synchronization} & Atomic instructions implement spin locks and higher-level primitives. Cache-line contention and memory ordering matter; coarse locks serialize the kernel, while fine locks increase complexity and deadlock risk.\\\hline
\textbf{8.1.4 Scheduling} & Time sharing allocates CPU intervals; space sharing grants processors for longer parallel phases; gang scheduling runs communicating threads together. Affinity preserves cache warmth but can resist load balance.\\\hline
\rowcolor{ice}\multicolumn{2}{l}{\hypertarget{sec-8-2}{}\textbf{Section 8.2: Multicomputers}}\\*\hline
\textbf{Private-memory model} & Each node owns processors and memory; nodes exchange messages over an interconnect. Explicit communication replaces shared loads and stores.\\\hline
\textbf{8.2.1 Hardware} & Interconnect topology, routing, switch design, network-interface processors, and DMA determine latency and throughput. Nodes may form grids, hypercubes, switched fabrics, or other networks.\\\hline
\textbf{8.2.2 Low-level software} & The kernel and interface coordinate send and receive buffers, descriptors, interrupts, DMA, retransmission, and flow control without racing over shared interface state.\\\hline
\textbf{8.2.3 User-level communication} & User libraries can bypass some kernel transitions by mapping network queues safely into a process. Protection and pinned buffers are needed to retain isolation and DMA correctness.\\\hline
\textbf{8.2.4 RPC} & A client stub marshals parameters into a request; the server stub unmarshals, calls the procedure, and returns results. Binding, representation, timeouts, duplicate requests, and partial failure complicate local-call transparency.\\\hline
\textbf{8.2.5 Distributed shared memory} & DSM maps shared pages or objects across nodes and uses faults plus messages to fetch, replicate, and invalidate data. False sharing and consistency protocol traffic can dominate.\\\hline
\textbf{8.2.6 Scheduling} & Placement must consider data, communication, processor capacity, and migration cost. Independent tasks can be distributed more freely than tightly communicating parallel tasks.\\\hline
\textbf{8.2.7 Load balancing} & Sender-initiated policies push work from overloaded nodes; receiver-initiated policies pull work when idle; bidding uses offers. Information freshness and transfer cost limit decisions.\\\hline
\rowcolor{ice}\multicolumn{2}{l}{\hypertarget{sec-8-3}{}\textbf{Section 8.3: Distributed Systems}}\\*\hline
\textbf{System character} & Independent computers cooperate over a network and fail partially. Middleware seeks a coherent service while naming, latency, concurrency, security, and failure remain distributed.\\\hline
\textbf{8.3.1 Network hardware} & Local and wide-area networks differ in topology, bandwidth, latency, error behavior, and administration. Packets cross links and routers rather than a shared memory fabric.\\\hline
\textbf{8.3.2 Services and protocols} & Layered protocols define message formats, sequencing, reliability, addressing, and transport. Connection-oriented and connectionless services offer different guarantees.\\\hline
\textbf{8.3.3 Document middleware} & Document-based systems use globally named resources, request/response protocols, and content formats; caching and replication improve reach but introduce freshness concerns.\\\hline
\textbf{8.3.4 File-system middleware} & A distributed file system gives remote files familiar operations. Client caching, stateless or stateful servers, naming, authentication, and consistency define observable semantics.\\\hline
\textbf{8.3.5 Object middleware} & Remote objects expose methods through proxies and brokers. Interface definitions, object references, marshaling, lifecycle, and exceptions support heterogeneous clients.\\\hline
\textbf{8.3.6 Coordination middleware} & A shared tuple or coordination space decouples producers and consumers by content rather than endpoint. Atomic insert, read, and remove operations coordinate distributed work.\\\hline
\rowcolor{ice}\multicolumn{2}{l}{\hypertarget{sec-8-4}{}\textbf{Section 8.4: Research on Multiple Processor Systems}}\\*\hline
\textbf{Research themes} & Research covers scalable coherence, weak memory, manycore scheduling, fast interconnects, one-sided communication, distributed consistency, failure handling, and energy-aware placement.\\\hline
\rowcolor{ice}\multicolumn{2}{l}{\hypertarget{sec-8-5}{}\textbf{Section 8.5: Summary}}\\*\hline
\textbf{Three scopes} & Multiprocessors share memory, multicomputers communicate explicitly between private memories, and distributed systems combine complete autonomous nodes through network services and middleware.\\\hline
\end{longtable}
\normalsize

\section*{Chapter Synthesis}
\begin{studybox}{Chapter Summary}
Multiple-processor systems range from tightly coupled shared-memory machines to private-memory multicomputers and geographically distributed nodes. As coupling weakens, ordinary loads and calls give way to explicit messages, RPC, DSM, and middleware. Across all forms, synchronization, placement, scheduling, load balance, consistency, and partial failure control scalability.
\end{studybox}

\begin{studybox}{Key Relationships}
\begin{itemize}
  \item Cache coherence makes shared memory usable, while synchronization defines legal concurrent updates.
  \item RPC builds a procedure abstraction over message passing but cannot remove network delay or partial failure.
  \item DSM builds a shared-memory abstraction over messages and therefore inherits both paging and consistency costs.
  \item Scheduling and data placement interact: moving computation may be cheaper than moving its working set.
\end{itemize}
\end{studybox}

\begin{studybox}{Design Tradeoffs}
\begin{itemize}
  \item Coarse locking is easier to reason about but limits multiprocessor concurrency.
  \item Processor affinity preserves locality, whereas migration improves balance.
  \item User-level networking cuts kernel overhead but requires careful protection and resource pinning.
  \item Transparent distributed abstractions simplify programming yet can hide latency and failure boundaries.
\end{itemize}
\end{studybox}

\begin{studybox}{Common Confusions}
\begin{itemize}
  \item A multiprocessor shares memory; a multicomputer has private memories connected by messages.
  \item Cache coherence concerns copies of memory blocks; consistency describes allowed visibility and ordering of operations.
  \item An RPC resembles a local call syntactically but has different failure, latency, and retry semantics.
\end{itemize}
\end{studybox}

\section*{Key Terms}
\small
\begin{longtable}{C N}
\rowcolor{navy}\color{white}\textbf{Term} & \color{white}\textbf{Concise definition}\\
\endfirsthead
\rowcolor{navy}\color{white}\textbf{Term} & \color{white}\textbf{Concise definition (continued)}\\
\endhead
\textbf{Multiprocessor} & Computer with multiple processors sharing a memory abstraction.\\\hline
\textbf{UMA} & Architecture with approximately uniform memory-access latency.\\\hline
\textbf{NUMA} & Architecture whose memory latency depends on physical placement.\\\hline
\textbf{Cache coherence} & Rules maintaining compatible values among cached copies.\\\hline
\textbf{SMP} & Organization in which processors run one shared operating system symmetrically.\\\hline
\textbf{Spin lock} & Lock acquired by repeatedly testing while remaining runnable.\\\hline
\textbf{Gang scheduling} & Concurrent scheduling of related threads on separate processors.\\\hline
\textbf{Multicomputer} & Private-memory nodes connected by a message network.\\\hline
\textbf{Marshaling} & Encoding parameters into a transferable representation.\\\hline
\textbf{RPC} & Request protocol presenting a remote operation as a procedure call.\\\hline
\textbf{DSM} & Shared-memory abstraction implemented across message-connected nodes.\\\hline
\textbf{False sharing} & Unrelated data causing coherence traffic because they share a transfer unit.\\\hline
\textbf{Load balancing} & Distribution of work to reduce overload and idle capacity.\\\hline
\textbf{Middleware} & Software layer offering distributed abstractions above host operating systems.\\\hline
\textbf{Partial failure} & Condition in which some distributed components fail while others continue.\\\hline
\end{longtable}
\normalsize

\enlargethispage{2\baselineskip}
\begin{studybox}{Active-Recall Review}
\small
\begin{enumerate}
  \item How do UMA and NUMA affect scheduling and memory placement?
  \item Why does an SMP kernel require synchronization?
  \item When should a waiting CPU spin rather than switch context?
  \item How does gang scheduling help communicating threads?
  \item What steps do RPC client and server stubs perform?
  \item Why can false sharing make DSM inefficient?
  \item How do sender- and receiver-initiated load balancing differ?
  \item Which transparency limits remain in distributed middleware?
  \item How do document, file, object, and coordination middleware differ?
\end{enumerate}
\end{studybox}

\par\medskip
{\color{navy}\bfseries Next-Chapter Connections}\par\smallskip
Chapter 9 focuses on protecting the processes, memory, files, devices, networks, and distributed interfaces introduced so far from accidental and malicious misuse.
\normalsize
\end{document}
